% ============================================================
% DRAFT — Appendix A: Project Planning Schedule.
%
% The PE448 2026 framework explicitly requires this ("Compile a
% planning schedule (Appendix A in project report) for the project"),
% and your own earlier status note already said this appendix is
% missing from your current draft. Two formats are both used in past
% Skripsie reports: a literal visual Gantt chart (e.g. DuPlessis 2020,
% just an inserted image/figure), or a week-by-week text table (e.g.
% Wait 2020, and the Theron 2025 report you just had me read). I've
% drafted the table version below because I have your real dates for
% it (from the six-week plan in the project's own status notes) and a
% table compiles with no extra packages. If you'd rather have an
% actual Gantt chart image instead, say so and I can build one
% (pgfgantt in LaTeX, or a plotted image you embed) from the same
% underlying dates -- just tell me which you want.
%
% Dates below come directly from your own project status notes
% (6-week shape, Sept-Oct 2026) and the PE448 2026 framework's fixed
% deadlines. The weekly goals are paraphrased from what actually
% happened/was planned, not invented -- but you know the ground truth
% better than I do, so check these against what you actually did,
% especially for the later weeks which were still "planned" rather
% than "done" as of the status note's date (16 Sept).
%
% This is a draft for review. Nothing here has been written into your
% actual chapter/appendix files.
% ============================================================

\chapter{Project Planning Schedule}
\label{app:planning}

The first meeting with my advisor took place at the start of the second
semester, where the project topic and scope were discussed and finalised.
Table~\ref{tab:planning} shows how the goals of the project were set out to
be completed.

\begin{table}[htbp]
    \centering
    \caption{Project Planning Schedule.}
    \label{tab:planning}
    \begin{tabular}{p{3.2cm}p{9.5cm}}
        \toprule
        Time Frame & Planned Goals \\
        \midrule
        16--22 Sep 2026 & Bench hardware sanity checks on the TurboPi: confirm
        full ROS2 bring-up, end-to-end motor control, and resolve the camera
        frame-rate issue caused by its auto-exposure behaviour in low light. \\
        \addlinespace
        23--29 Sep 2026 & Finish the physical rover controller (drive and
        camera), wire in line-following, and get the first full
        voice-to-drive-to-report loop working on the simplest command level. \\
        \addlinespace
        30 Sep--6 Oct 2026 & Harden the physical integration (line lost/
        regained, ultrasonic obstacle stop, multi-waypoint routes); draft the
        Implementation chapter. \\
        \addlinespace
        7--13 Oct 2026 & Run the full experiment battery: the physical command
        levels (L0--L5) and the virtual compositional-depth sweep; generate
        trace and result figures. \\
        \addlinespace
        14--20 Oct 2026 & Write the Evaluation and Conclusion chapters;
        consistency pass on the Introduction and Literature Review. \\
        \addlinespace
        21--27 Oct 2026 & Full report compile and proofreading pass; draft the
        poster and oral presentation outline. \\
        \addlinespace
        28--30 Oct 2026 & Final editing and formatting pass; report submission
        (hand-in deadline 30 October 2026, 23:59). \\
        \bottomrule
    \end{tabular}
\end{table}

Oral presentations and the internal examination take place 6--13 November
2026, followed by the project open day (poster presentation and external
moderator interviews) on 18 November 2026.
